\documentclass[12pt]{article}
\usepackage[utf8]{inputenc}
\usepackage[T1]{fontenc}
\usepackage{hyperref}
\title{Tira Language Resources for K-12 Curriculum Development}
\author{Compiled by Claude Sonnet 5 (claude-sonnet-5) for Andr\'as Kornai}
\date{}
\begin{document}
\maketitle
\section{Introduction}
Tira (also spelled Tiro or Thiro; ISO 639-3 tic) is a Heiban-family language of the Nuba Mountains, Sudan. Wikipedia gives about 34,000 speakers (2022) and a Latin-script writing system. The search found no Tira-specific school materials, so this dossier is short.
\section{Primary Educational Resources}
No Tira primers, textbooks or literacy programmes turned up in the search results.
\section{Linguistic and Reference Materials}
\begin{itemize}
\item R. C. Stevenson, \emph{Tira and Otoro: Two Kordofanian Grammars}, edited by Thilo Schadeberg. It is based on 1942--43 fieldwork and includes example texts and short glossaries. It is listed in the \emph{Nuba Mountain Language Studies} series: \url{https://koeppe.de/titel_nuba-mountain-language-studies}
\item \emph{Nuba Mountain Language Studies} (Schadeberg and Blench, eds., 2013), a collection of 32 papers. Its Heiban section includes a paper on Moro literacy, which is a different language and is relevant only as a comparison: \url{https://www.koeppe.de/titel_print_nuba-mountain-language-studies}
\item Wikipedia overview of Tiro: \url{https://en.wikipedia.org/wiki/Tiro_language}
\end{itemize}
\section{Religious and Community Resources}
Joshua Project has a people-group profile for the Tira in Sudan: \url{https://joshuaproject.net/people_groups/15499}. No scripture translations or community organisations were found.
\section{Digital and Audio Resources}
None found.
\section{Curriculum Development Context}
Many Tira speakers also use Sudanese Arabic. No evidence was found of Tira being taught in schools. A curriculum would have to be built from scratch, starting from the Stevenson grammar and new fieldwork.
\end{document}